% ============================================================
%  SCS 3312 – Research Methods  |  D0: Topic Proposal & Research Plan
%  Group Submission
% ============================================================
\documentclass[12pt, a4paper]{report}

% ── Packages ────────────────────────────────────────────────
\usepackage[left=1.5in, right=1in, top=1.5in, bottom=1in]{geometry}
\usepackage{fontenc}
\usepackage{inputenc}
\usepackage{setspace}
\usepackage{titlesec}
\usepackage{tocloft}
\usepackage{fancyhdr}
\usepackage{hyperref}
\usepackage{bookmark}
\usepackage{parskip}
\usepackage{enumitem}
\usepackage{booktabs}
\usepackage{array}
\usepackage{makeidx}
\usepackage{natbib}   % Harvard-style citations

\makeindex

% ── Hyperref / Bookmarks ────────────────────────────────────
\hypersetup{
    colorlinks = true,
    linkcolor  = black,
    citecolor  = black,
    urlcolor   = blue,
    bookmarks  = true,
    bookmarksopen = true,
    pdftitle   = {Passenger Waiting Behaviour and Boarding Decisions at Selected Urban Bus Stops},
    pdfauthor  = {Group – SCS 3312},
}

% ── Page numbering style helper ─────────────────────────────
% Roman numerals for front matter, Arabic from Chapter 1 onwards
% (handled via \frontmatter / \mainmatter in the document body)

% ── Header / Footer ─────────────────────────────────────────
\pagestyle{fancy}
\fancyhf{}
\renewcommand{\headrulewidth}{0pt}
\cfoot{\thepage}          % page number centred in footer

% Title page has no page number
\fancypagestyle{plain}{
    \fancyhf{}
    \renewcommand{\headrulewidth}{0pt}
    \cfoot{}
}

% ── Chapter / Section title formatting ──────────────────────
\titleformat{\chapter}[block]
  {\normalfont\Large\bfseries}{\chaptername\ \thechapter}{1em}{}
\titlespacing*{\chapter}{0pt}{30pt}{20pt}

% ── Line spacing ─────────────────────────────────────────────
\onehalfspacing

% ============================================================
\begin{document}

% ────────────────────────────────────────────────────────────
%  TITLE PAGE  (no page number)
% ────────────────────────────────────────────────────────────
\begin{titlepage}
    \centering
    \vspace*{2cm}

    {\Large\bfseries University of Colombo School of Computing\\[0.4em]
    SCS 3312 – Research Methods}

    \vspace{2cm}

    \rule{\linewidth}{0.5pt}\\[0.6em]
    {\LARGE\bfseries Passenger Waiting Behaviour and\\[0.3em]
    Boarding Decisions at Selected\\[0.3em]
    Urban Bus Stops}\\[0.4em]
    \rule{\linewidth}{0.5pt}

    \vspace{1.5cm}

    {\large\textbf{D0: Topic Proposal \& Research Plan}}\\[0.3em]
    {\large Group Submission}

    \vspace{2cm}

    \begin{tabular}{ll}
        \textbf{Module}     & SCS 3312 – Research Methods \\[0.3em]
        \textbf{Deadline}   & 20 May 2026, 12:00 Hrs \\[0.3em]
        \textbf{Submission} & Group \\
    \end{tabular}

    \vfill

    {\large University of Colombo School of Computing\\
    University of Colombo, Sri Lanka}\\[0.5em]
    {\large May 2026}
\end{titlepage}

% ────────────────────────────────────────────────────────────
%  FRONT MATTER – Roman numerals, no "Chapter" prefix
% ────────────────────────────────────────────────────────────
\pagenumbering{roman}
\setcounter{page}{1}

% No section numbering before Chapter 1
\setcounter{secnumdepth}{-1}

% ── Acknowledgement ─────────────────────────────────────────
\chapter*{Acknowledgement}
\addcontentsline{toc}{chapter}{Acknowledgement}

We would like to express our sincere gratitude to the faculty of the University of Colombo
School of Computing for providing this opportunity to engage with real-world research
environments. We also thank the passengers and commuters at the selected bus stops whose
cooperation made field observations possible, as well as the University of Colombo
administration for facilitating access to the study sites.

% ── Abstract ────────────────────────────────────────────────
\chapter*{Abstract}
\addcontentsline{toc}{chapter}{Abstract}

This proposal presents a research plan to investigate passenger waiting behaviour and
boarding decisions at three selected urban bus stops in Colombo, Sri Lanka. The study is
motivated by the observable yet under-studied phenomena of passengers extending their
waiting duration, rejecting overcrowded buses, and exhibiting irregular boarding patterns
within public transport environments. Using a mixed-method approach that combines
qualitative techniques—direct observation and informal interviews—with quantitative
methods—passenger counting, waiting-time measurement, and structured surveys—the
research aims to identify behavioural patterns and the factors that influence them. Three
field sites, namely Reid Avenue, Town Hall, and Fort Main bus stops, have been selected
because they represent varied levels of passenger congestion and operational conditions.
Data will be collected during both peak and off-peak hours to capture the full range of
commuter behaviour. The findings are expected to contribute practical insights for improving
passenger management and public transport planning in urban Sri Lanka.

% ── Table of Contents ───────────────────────────────────────
\tableofcontents

% ── List of Figures ─────────────────────────────────────────
\listoffigures
\addcontentsline{toc}{chapter}{List of Figures}

% ── List of Tables ──────────────────────────────────────────
\listoftables
\addcontentsline{toc}{chapter}{List of Tables}

% ── List of Acronyms ────────────────────────────────────────
\chapter*{List of Acronyms}
\addcontentsline{toc}{chapter}{List of Acronyms}

\begin{tabular}{ll}
    \textbf{UCSC} & University of Colombo School of Computing \\
    \textbf{LMS}  & Learning Management System \\
    \textbf{SCS}  & School of Computing Studies \\
\end{tabular}

% ────────────────────────────────────────────────────────────
%  MAIN MATTER – Arabic numerals, section numbering ON
% ────────────────────────────────────────────────────────────
\pagenumbering{arabic}
\setcounter{page}{1}
\setcounter{secnumdepth}{3}   % number sections, subsections, subsubsections

% ============================================================
%  CHAPTER 1 – INTRODUCTION
% ============================================================
\chapter{Introduction}

This chapter provides an overview of the research study and introduces the main problem
being investigated. It explains the background and importance of studying passenger waiting
behaviour at urban bus stops, presents the research objectives and questions, and defines
the scope of the study. Furthermore, this chapter establishes the conceptual foundation for
understanding how passenger behaviour is shaped by factors such as waiting time, crowd
conditions, and bus availability in real-world transport environments.

\section{Background of the Study}

Public transportation is one of the most essential services used by people in urban areas.
In Sri Lanka, buses serve as the primary mode of daily travel for students, office workers,
and the general public. Consequently, bus stops have become crowded public spaces where
passengers wait, interact, and make consequential decisions relating to transportation.

Passenger behaviour at bus stops varies considerably depending on conditions such as
waiting time, crowd level, bus availability, and urgency of travel. Some passengers wait
for a specific bus route for extended periods, while others board the first available bus. In
crowded situations, passengers may avoid overcrowded buses, move frequently within the
waiting area, or attempt aggressive boarding when buses arrive.

These behaviours have direct implications for passenger safety, comfort, queue discipline,
and the overall efficiency of public transportation services. Studying passenger waiting
behaviour therefore offers a window into how people make real-world transport decisions
under varying environmental pressures.

This study focuses on observing and analysing passenger waiting behaviour and boarding
decisions at three selected urban bus stops in Colombo during both peak and off-peak
hours.

\section{Research Problem}

Passenger behaviour at bus stops is frequently influenced by overcrowding, irregular bus
arrivals, prolonged waiting times, and individual travel preferences. Although buses may
arrive with reasonable frequency, many passengers continue to wait for specific services
rather than boarding immediately. Some passengers reject visibly overcrowded buses, while
others display visible impatience during extended waits.

At congested bus stops, passengers may also demonstrate behaviours such as queue
changing, standing at the roadside, relocating between waiting areas, or rushing to board
arriving buses. These behaviours can create safety hazards, increase congestion, and
reduce boarding efficiency.

Despite the prevalence of these phenomena, there is limited scholarly understanding of the
factors that govern passenger waiting decisions and behavioural responses in urban bus
stop environments. This research therefore aims to investigate why passengers extend their
waiting periods and how different environmental and social factors affect their behaviour at
bus stops.

\section{Research Topic}

\begin{center}
    \textit{"Passenger Waiting Behaviour and Boarding Decisions\\
    at Selected Urban Bus Stops"}
\end{center}

This research examines how passengers behave while waiting for buses, and how factors
such as crowd density, waiting duration, and bus conditions influence their boarding
decisions.

\section{Relevance of the Study}

Public buses constitute one of the primary transportation modes used in Colombo. Passenger
behaviour at bus stops directly affects crowd management, transport efficiency, and passenger
safety. Understanding waiting behaviour can help to:

\begin{itemize}[noitemsep]
    \item Identify common behavioural patterns at bus stops;
    \item Understand the factors that cause prolonged waiting times;
    \item Improve passenger management strategies;
    \item Support better public transport planning; and
    \item Reduce overcrowding and unsafe boarding behaviour.
\end{itemize}

In addition, this research is relevant to the field of research methods because it studies
observable, real-world human behaviour using both qualitative and quantitative approaches,
making it a suitable vehicle for developing methodological competence.

\section{Research Objectives}

The primary objective of this study is to investigate passenger waiting behaviour and boarding
decisions at selected bus stops in Colombo. The specific objectives are:

\begin{enumerate}[noitemsep]
    \item To identify common passenger waiting behaviours at bus stops;
    \item To examine the factors influencing passenger waiting time;
    \item To analyse how crowd density affects boarding decisions;
    \item To compare passenger behaviour between different bus stops; and
    \item To observe behavioural differences during peak and off-peak hours.
\end{enumerate}

\section{Preliminary Research Questions}

This study attempts to address the following research questions:

\begin{enumerate}[noitemsep]
    \item Why do some passengers wait longer than others at bus stops?
    \item How does crowd density influence passenger boarding decisions?
    \item Why do some passengers reject overcrowded buses?
    \item How do passengers behave during extended waiting periods?
    \item Are there observable differences in passenger behaviour between different bus stops?
    \item How do peak-hour conditions affect passenger waiting behaviour?
\end{enumerate}

\section{Scope of the Study}

This research is limited to observing passenger waiting behaviour and boarding decisions at
three selected urban bus stops in Colombo. The study focuses on:

\begin{itemize}[noitemsep]
    \item Passenger waiting time;
    \item Boarding and rejection behaviour;
    \item Crowd conditions within waiting areas;
    \item Passenger movement within bus stops; and
    \item Peak-hour and off-peak-hour behaviour.
\end{itemize}

Data will be collected through direct observation, informal interviews, surveys, passenger
counting, and waiting-time measurements. The study does not encompass technical bus
operations, government transport policies, or driver behaviour. The primary focus remains
on passenger behaviour and decision-making within bus stop environments.

% ============================================================
%  CHAPTER 2 – FIELD STUDY DESCRIPTION
% ============================================================
\chapter{Field Study Description}

This chapter describes the locations selected for the research study and explains the
environmental conditions observed at each bus stop. It discusses the differences between
peak-hour and off-peak-hour conditions, and identifies common passenger activities and
behavioural patterns noted during preliminary field visits. The purpose of this chapter is to
provide a clear understanding of the real-world setting in which the research will be
conducted.

\section{Selected Field Sites}

The research will be conducted at three urban bus stops in Colombo. These sites were
selected because they experience different passenger volumes, bus frequencies, and crowd
conditions throughout the day, thereby providing suitable environments for observing a
range of passenger behaviours related to waiting, boarding, and decision-making.

The selected field sites are:

\begin{itemize}[noitemsep]
    \item Reid Avenue Bus Stop
    \item Town Hall Bus Stop
    \item Fort Main Bus Stop
\end{itemize}

These locations are commonly used by students, office workers, and the general public. The
bus stops also represent different levels of passenger congestion, allowing meaningful
comparisons between highly crowded and moderately crowded environments.

\section{Description of the Selected Bus Stops}

Each selected bus stop has distinct physical and environmental characteristics that may
influence passenger behaviour.

\subsection{Reid Avenue Bus Stop}

This bus stop is primarily used by university students and nearby office workers. Passenger
traffic increases significantly during morning and evening hours. The waiting space is
limited, which sometimes causes crowding near the roadside and contributes to disorganised
boarding behaviour.

\subsection{Town Hall Bus Stop}

Town Hall Bus Stop is situated in a busy commercial and administrative area. A large number
of passengers use this location throughout the day. Due to the high number of bus routes
served, passengers frequently move within the waiting area and often change their boarding
decisions in response to crowd conditions.

\subsection{Fort Main Bus Stop}

Fort Main Bus Stop experiences very high passenger flow as it functions as a major transport
hub in Colombo. Passengers at this location regularly encounter long waiting times,
overcrowding, and competition for boarding space during peak hours.

These differences between bus stops provide opportunities to compare passenger behaviour
under contrasting environmental and operational conditions.

\section{Peak and Off-Peak Conditions}

Passenger behaviour at bus stops changes substantially depending on the time of day and
prevailing crowd conditions. The study will therefore observe both peak-hour and
off-peak-hour situations.

\subsection{Peak Hours}

Peak hours typically occur during morning travel periods and evening return hours. During
these periods:

\begin{itemize}[noitemsep]
    \item Passenger density increases markedly;
    \item Waiting areas become congested;
    \item Boarding competition intensifies; and
    \item Passengers may become impatient as a result of delays or overcrowding.
\end{itemize}

\subsection{Off-Peak Hours}

Off-peak periods generally exhibit lower passenger volumes, shorter waiting times, reduced
crowd pressure, and more organised boarding behaviour. Observing both conditions will help
identify how environmental factors influence waiting behaviour and decision-making.

\section{Observed Passenger Activities}

During preliminary observations, several common passenger activities and behavioural
patterns were identified at the selected bus stops. These include:

\begin{itemize}[noitemsep]
    \item Waiting for specific bus routes;
    \item Rejecting overcrowded buses upon arrival;
    \item Queue changing and position shifting;
    \item Moving closer to the roadside when buses approach;
    \item Using mobile phones while waiting;
    \item Interacting with fellow passengers;
    \item Attempting aggressive boarding during congested situations; and
    \item Frequently checking for approaching buses.
\end{itemize}

Differences were also observed between passengers who were prepared to wait for extended
periods and those who preferred immediate boarding regardless of conditions. These
behaviours indicate that passenger decisions are shaped by waiting time, crowd conditions,
comfort, and urgency of travel, and they establish a firm foundation for the qualitative and
quantitative data collection planned for subsequent phases of the study.

% ============================================================
%  CHAPTER 3 – RESEARCH METHODOLOGY
% ============================================================
\chapter{Research Methodology}

This chapter explains the research methods and procedures that will be used to collect and
analyse data. It describes both qualitative and quantitative approaches to investigating
passenger waiting behaviour and boarding decisions at selected bus stops, and outlines
methods of observation, interviews, surveys, data recording, and the responsibilities of
each group member.

\section{Research Approach}

This research adopts a mixed-method approach, combining qualitative and quantitative data
collection. The mixed-method design was selected because passenger behaviour at bus
stops involves both observable behavioural patterns and measurable numerical data
\citep{creswell2014}.

The qualitative strand will illuminate passenger experiences, behaviours, and
decision-making processes through direct observation and informal interviews. The
quantitative strand will produce measurable data on waiting times, passenger counts, and
boarding frequencies. By integrating both strands, the study aims to achieve a more
comprehensive understanding of passenger waiting behaviour under varying transport
conditions.

\section{Qualitative Data Collection Methods}

The qualitative phase focuses on understanding passenger behaviour, attitudes, and
decision-making patterns. Data will be collected through observation, informal interviews,
and field notes.

\subsection{Observation}

Direct observation will be used to study passenger activities and behavioural patterns at the
selected bus stops during both peak and off-peak hours. Observations will focus on:

\begin{itemize}[noitemsep]
    \item Passenger movement within the bus stop;
    \item Queue behaviour and queue discipline;
    \item Boarding behaviour;
    \item Bus rejection behaviour;
    \item Reactions to delays and overcrowding; and
    \item Passenger interactions with others.
\end{itemize}

Observation allows researchers to study natural passenger behaviour in real-world settings
without interfering with normal activities \citep{bryman2016}.

\subsection{Informal Interviews}

Short informal interviews will be conducted with selected passengers. These interviews will
help to clarify the reasons behind observed decisions and behaviours. Topics will include:

\begin{itemize}[noitemsep]
    \item Reasons for waiting for a specific bus;
    \item Opinions regarding overcrowded buses;
    \item Factors affecting boarding decisions; and
    \item Passenger experiences during long waiting periods.
\end{itemize}

Interviews will be brief, respectful, and conducted without disturbing passengers.

\subsection{Field Notes}

Field notes will record important observations and environmental conditions during data
collection, including crowd conditions, weather, unusual passenger behaviours, traffic
conditions, and any special incidents observed. Field notes will support the interpretation of
qualitative findings in later analytical stages.

\section{Quantitative Data Collection Methods}

The quantitative phase focuses on collecting measurable numerical data related to
passenger waiting behaviour.

\subsection{Passenger Counting}

Passenger counting will measure the number of passengers present at each bus stop across
different time periods, including the number arriving, boarding, and rejecting buses, as well
as passenger density during peak and off-peak hours. This data will enable identification of
crowd patterns and passenger flow.

\subsection{Waiting Time Measurement}

Waiting-time measurements will determine how long passengers remain at bus stops before
boarding. Records will capture average and maximum waiting times, differences between
peak and off-peak periods, and waiting durations for different bus routes.

\subsection{Survey Data Collection}

Structured questionnaires will be distributed to selected passengers. Survey questions will
address average waiting time, preferred bus conditions, willingness to board crowded buses,
passenger satisfaction, and reasons for rejecting buses. Responses will be summarised
using percentages, charts, and basic statistical comparisons.

\section{Variables to be Measured}

The study will measure the following variables:

\begin{itemize}[noitemsep]
    \item Passenger waiting time;
    \item Passenger count;
    \item Crowd density;
    \item Bus arrival intervals;
    \item Number of rejected buses;
    \item Boarding frequency; and
    \item Peak-hour and off-peak-hour passenger flow.
\end{itemize}

\section{Data Recording Procedure}

Data collection will be conducted systematically at selected bus stops at different times of
day through the following procedure:

\begin{enumerate}[noitemsep]
    \item Visiting selected bus stops during peak and off-peak hours;
    \item Observing passenger behaviour directly;
    \item Recording passenger counts and waiting times;
    \item Conducting short interviews and distributing surveys;
    \item Maintaining field notes for each observation session; and
    \item Organising collected data into tables and records for analysis.
\end{enumerate}

Multiple observation visits will be conducted at each location to improve accuracy and
reliability.

\section{Group Member Responsibilities}

Research tasks will be distributed equitably among the three group members to ensure
efficient data collection and coordination.

\begin{description}
    \item[Member 1 – Observation and Passenger Counting]
        Observe passenger behaviour; count passengers during selected time periods;
        measure crowd density; and record boarding and rejection behaviour.

    \item[Member 2 – Interviews and Survey Collection]
        Conduct informal interviews; distribute and collect survey questionnaires;
        record passenger opinions and responses; and organise qualitative data.

    \item[Member 3 – Waiting Time and Field Documentation]
        Measure passenger waiting times; record bus arrival intervals; maintain field notes;
        and document environmental conditions and special observations.
\end{description}

All group members will participate in field visits and collaborate during data analysis and
report preparation.

% ============================================================
%  CHAPTER 4 – CRITICAL ANALYSIS
% ============================================================
\chapter{Critical Analysis}

This chapter critically discusses the importance of studying passenger waiting behaviour and
evaluates the practical challenges and limitations associated with the research. It also
highlights the ethical considerations that must be followed during data collection.

\section{Importance of Studying Passenger Behaviour}

Public transportation is central to urban daily life, and passenger behaviour at bus stops
directly affects crowd management, boarding efficiency, passenger safety, and the overall
quality of the public transport experience. Studying waiting behaviour can reveal how people
respond to overcrowding, long waiting times, and irregular bus arrivals, thereby identifying
systemic transport problems.

This study is important because it:

\begin{itemize}[noitemsep]
    \item Elucidates real-world commuter decision-making processes;
    \item Identifies behavioural patterns under crowded conditions;
    \item Provides insights into passenger comfort and safety concerns;
    \item Contributes to improving public transportation experiences; and
    \item Supports the development of more efficient passenger management strategies.
\end{itemize}

From an academic perspective, the research is valuable because it applies both qualitative
and quantitative methods to study real-world human behaviour in public environments,
aligning with contemporary mixed-method research traditions \citep{tashakkori2010}.

\section{Expected Challenges During Data Collection}

Several challenges are anticipated. Overcrowding during peak hours may make it difficult to
observe behaviours accurately, count passengers, and conduct interviews without
disruption. The unpredictable nature of bus transportation—including delays, traffic
congestion, and sudden changes in passenger flow—may affect the consistency of
observations and waiting-time measurements.

Passengers may also be unwilling to participate in interviews or surveys due to time
constraints, privacy concerns, or lack of interest, and some may provide incomplete or
inaccurate responses. Environmental conditions such as rain, noise, heat, and traffic
congestion may hinder field observations and data recording. Managing multiple concurrent
activities—observation, counting, and interviews—will require careful coordination among
group members.

\section{Possible Limitations of the Study}

The research is limited to three selected bus stops in Colombo; the observed behavioural
patterns may therefore not fully represent passenger behaviour across all bus stops or
transport environments in Sri Lanka. The observation period is restricted to selected days
and times, meaning that behaviour during weekends, holidays, special events, or atypical
transport conditions falls outside the study scope.

Additionally, certain passenger decisions and emotional responses cannot be directly
measured through observation alone; some behaviours may be influenced by personal
factors that remain invisible during fieldwork. The accuracy of waiting-time measurements
and passenger counts may also be affected by heavy crowd movement and rapid changes in
passenger flow during peak hours. Despite these limitations, the study is expected to yield
meaningful insights into passenger waiting behaviour and boarding decisions in urban
transport environments.

\section{Ethical Considerations}

The research will adhere to ethical standards throughout all stages of data collection.
Passenger privacy and confidentiality will be respected; no personal information such as
names, contact details, or identification will be collected. Informal interviews and surveys
will be conducted only with the informed consent of participants, who retain the right to
withdraw at any time.

Researchers will avoid disrupting normal passenger activities during observations, which
will focus exclusively on general public behaviour in open, publicly accessible spaces. All
collected data will be used solely for the academic purposes of this study, and findings will
be reported honestly and accurately without manipulation or misrepresentation.

\end{document}